\question[6]
Vergleiche die folgenden Verschlüsselungsverfahren:

\begin{itemize}
\item Transpositionsverfahren (z.\,B. Skytale)
\item monoalphabetische Substitution (z.\,B. Caesar-Verfahren)
\item polyalphabetische Substitution (z.\,B. Vigenère-Verfahren)
\end{itemize}

Erläutere jeweils, wie die Verschlüsselung erfolgt, und nenne einen Vor- bzw. Nachteil des Verfahrens.

\begin{solutionbox}{6cm}

Transposition: Die Buchstaben des Klartextes bleiben erhalten, ihre Reihenfolge wird verändert. \\
Monoalphabetische Substitution: Jeder Buchstabe wird im gesamten Text durch denselben Ersatzbuchstaben ersetzt. \\
Polyalphabetische Substitution: Die Ersetzung eines Buchstabens hängt von seiner Position bzw. vom Schlüssel ab.  

Mögliche Vor- und Nachteile: \\
Transposition: einfache Durchführung, aber häufig leicht zu entschlüsseln. \\
Monoalphabetische Substitution: einfach anzuwenden, aber anfällig für Häufigkeitsanalysen. \\
Polyalphabetische Substitution: sicherer gegen Häufigkeitsanalysen, jedoch aufwendiger.
\end{solutionbox}